\chapter{Literature Review and Related Work}
\label{chap:lit_review}

\section{Review Approach}
To establish a rigorous theoretical and empirical foundation for the CycleSafe platform, a systematic literature survey was conducted across peer-reviewed transportation science journals, intelligent transportation systems (ITS) conferences, and open-source hardware repositories. The review focused on four primary literature thematic clusters published between 2010 and 2026:
\begin{enumerate}
    \item Naturalistic cycling studies measuring physical motor vehicle overtaking distances.
    \item Embedded IoT proximity sensing nodes and hardware telemetry units deployed on non-motorized vehicles.
    \item Smartphone-based crowdsourcing and inertial near-miss detection platforms.
    \item Microscopic multi-modal traffic simulation frameworks modeling bicycle-vehicle interactions.
\end{enumerate}

\section{Thematic Review}

\subsection{Naturalistic Overtaking Studies and Passing Distance Laws}
Over the past decade, empirical naturalistic cycling studies have established that motorized overtaking clearance directly governs cyclist perceived safety and objective injury risk \cite{chataway2014passing, dozza2016what}. Early pioneering research by Walker utilized instrumented bicycles equipped with ultrasonic rangefinders and video cameras to quantify driver behavior, revealing that vehicle speed, oncoming traffic density, and lane width significantly constrict overtaking margins. 

Following the global proliferation of Minimum Passing Distance Laws (MPDL)---exemplified by Germany's StVO \S\,5(4) mandating a 1.5\,m clearance in urban zones \cite{stvo2020germany} and Australia's 1.0\,m/1.5\,m statutory rules \cite{beck2019passing}---researchers turned to dedicated telemetry devices to assess compliance. Beck et al.\ \cite{beck2019passing} demonstrated in a multi-city Australian naturalistic cycling cohort that between 10\% and 19\% of all overtaking events violated statutory passing distance laws, with higher violation rates occurring on roads without marked bicycle lanes and along curved segments.

\subsection{Embedded Hardware Sensing: The OpenBikeSensor (OBS) Platform}
The most prominent open-hardware initiative targeting passing distance measurement is the \textbf{OpenBikeSensor (OBS)} project \cite{ehrenwirth2020openbikesensor}. Developed by civic technology communities in Germany, OBS employs an Espressif ESP32 dual-core microcontroller interfaced with two lateral waterproof ultrasonic transducers (JSN-SR04T), a u-blox satellite receiver, a small OLED display, and a handlebar momentary push-button.

In the OBS workflow, ultrasonic pulses continuously measure lateral clearance to the left (overtaking motorized vehicles) and to the right (parked vehicles and roadside infrastructure). When an overtaking event occurs, the rider actuates the handlebar button, causing the firmware to record a confirmation marker in the telemetry log. While OBS represents a milestone in citizen science and community data collection, several critical engineering limitations remain unaddressed in the literature:
\begin{itemize}
    \item \textbf{Absence of Inertial Tilt Compensation:} Bicycles lean into turns and oscillate laterally during active pedaling. When a bicycle tilts at a roll angle $\phi$, an ultrasonic transducer rigidly attached to the frame emits an acoustic cone directed obliquely downward toward the pavement or upward toward the vehicle roofline. Without real-time accelerometer roll compensation, this geometric tilt introduces systematic distance measurement errors.
    \item \textbf{Power Rail Instability and SD Write Corruption:} Operating the ESP32 Wi-Fi radio, ultrasonic transducers, and MicroSD card SPI bus from unbuffered battery rails frequently induces transient brownout resets, causing FAT32 file allocation table corruptions during long recording sessions.
\end{itemize}

\subsection{Crowdsourced Smartphone Incident Sensing: The SimRA Platform}
As an alternative to specialized hardware sensors, researchers at Technische Universit{\"a}t Berlin developed \textbf{SimRA} (Safety in Mobility: Ride Analytics) \cite{simra2020crowdsourcing}. SimRA utilizes the embedded tri-axial accelerometers and GPS receivers of consumer smartphones mounted on bicycle handlebars to passively detect sudden braking ($\Delta a_x < -2.5\,\text{m/s}^2$) and sharp swerving maneuvers ($\Delta a_y$). At the conclusion of a ride, the cyclist annotates flagged near-miss incidents with contextual details (e.g., cut-ins, car door openings, close passes).

While SimRA successfully engaged tens of thousands of participants and collected millions of kilometers of urban bicycle trajectory data across European cities, its reliance on smartphones incurs fundamental physical limitations:
\begin{enumerate}
    \item Smartphones possess no active ranging sensors (ultrasonic or radar); consequently, passing distance cannot be measured physically and relies entirely on subjective post-ride estimation.
    \item Smartphone mounting orientations vary widely, and mechanical phone holders introduce high-frequency vibrational noise that complicates deterministic acceleration thresholding.
\end{enumerate}

\subsection{Microscopic Multi-Modal Traffic Simulation in SUMO}
Microscopic traffic simulation offers a reproducible computational laboratory for modeling complex vehicular interactions without the logistical overhead and safety risks of on-road trials. \textbf{Eclipse SUMO (Simulation of Urban MObility)} \cite{krajzewicz2012recent, lopez2018microscopic} is the leading open-source microscopic simulation suite. SUMO implements car-following models such as the Intelligent Driver Model (IDM) \cite{treiber2000congested} and the Krau{\ss} model, alongside the Sublane Model, which divides discrete road lanes into continuous lateral strips (e.g., 0.2\,m resolution) to simulate realistic non-lane-based micro-mobility overtaking.

Through the Traffic Control Interface (TraCI), external control routines can subscribe to moving vehicle variables (spatial coordinates, instantaneous speeds, acceleration vectors, and contextual distance envelopes) at each simulation step. However, existing SUMO literature predominantly focuses on motorized vehicle delay minimization and intersection throughput, rarely modeling the continuous lateral acoustic clearance and inertial roll dynamics experienced by individual cyclists in mixed-traffic streams.

\section{Comparative Analysis of Related Approaches}
To contextualize the CycleSafe contribution within the current technological landscape, Table \ref{tab:lit_comparison} presents a comprehensive comparative benchmark against existing platforms.

\begin{table}[h!]
	\centering
	\caption{Comparative benchmarking of cycling safety evaluation platforms.}
	\label{tab:lit_comparison}
	\setlength{\tabcolsep}{6pt}
	\renewcommand{\arraystretch}{1.3}
	\begin{adjustbox}{width=\textwidth}
	\begin{tabular}{|p{3.2cm}|p{2.6cm}|p{2.8cm}|p{2.6cm}|p{3.4cm}|}
		\hline
		\textbf{Evaluation Dimension} & \textbf{ProxiCycle} & \textbf{SimRA (TU Berlin)} & \textbf{OpenBikeSensor (OBS)} & \textbf{CycleSafe (This Work)} \\
		\hline
		\hline
		\textbf{Core Hardware Architecture} & Arduino / Raspberry Pi & Smartphone (Android/iOS) & ESP32 Dual-Core DevKit & ESP32 + MPU-6050 6-Axis IMU \\
		\hline
		\textbf{Ranging Technology} & Infrared / Ultrasonic & None (Inertial only) & Dual Ultrasonic (JSN-SR04T) & Dual Ultrasonic (JSN-SR04T) \\
		\hline
		\textbf{Distance Output} & Binary (Passing / Clear) & Subjective rider annotation & Continuous metric distance (cm) & Continuous metric distance (cm) \\
		\hline
		\textbf{Dynamic Tilt Compensation} & None & None & None & Real-time Accelerometer Roll Model ($\phi = \arctan(a_y / \sqrt{a_x^2 + a_z^2})$) \\
		\hline
		\textbf{Vehicle Class Discrimination} & None & User post-ride survey & Generic ultrasonic envelope & Supervised XGBoost Classifier (Truck / Car / Moto / Bike) \\
		\hline
		\textbf{Microscopic Simulation Twin} & None & None & None & 10\,km SUMO TraCI Digital Twin (1,244 Vehicles) \\
		\hline
		\textbf{Power Supply Reliability} & Unregulated 5\,V Rail & Phone internal battery & Shared single rail & Dedicated AMS1117-3.3V LDO Rail \\
		\hline
		\textbf{Web Telemetry Export} & None & Web portal (KML/CSV) & Web portal (JSON) & Full-Stack REST API + Instant CSV Export \\
		\hline
	\end{tabular}
	\end{adjustbox}
\end{table}

\section{Identified Research Gaps}
The literature survey identifies three primary research gaps that directly motivate this investigation:
\begin{enumerate}
    \item \textbf{Gap 1: Geometric Tilt Distortion in Bicycle Ranging Telemetry.} Existing hardware proximity platforms (including OpenBikeSensor and ProxiCycle) assume the instrumented bicycle operates perpetually in an upright perpendicular posture. In real-world urban cycling, cornering, maneuvering around road obstacles, and standard pedaling cadence produce dynamic roll angles between $5^\circ$ and $25^\circ$, introducing systematic distance overestimation that invalidates close-pass statistical baselines.
    \item \textbf{Gap 2: Uniform Severity Treatment Across Disparate Vehicle Classes.} Conventional threshold-based collision warning systems treat all lateral clearances below 1.5\,m identically. However, in urban freight corridors, an overtaking heavy-duty truck passing at 1.4\,m presents vastly greater turbulence, kinetic mass threat, and cognitive panic than a peer bicycle passing at 1.4\,m on a separated path. There is a lack of machine learning models that fuse physical clearance with counterpart vehicle taxonomy.
    \item \textbf{Gap 3: Decoupling of Empirical IoT Logging and Microscopic Simulation.} Naturalistic IoT telemetry collection and microscopic simulation modeling have progressed along isolated tracks. Physical sensor data is rarely used to parameterize and validate multi-modal simulation twins, and traffic simulation platforms rarely generate synthetic telemetry matching the exact schema and physics of embedded hardware nodes.
\end{enumerate}

\section{Chapter Summary}
This chapter reviewed the state of the art in cyclist overtaking safety evaluation, spanning statutory passing regulations, embedded telemetry hardware, smartphone crowdsourcing, and microscopic traffic simulation. The critical shortcomings of existing solutions---lack of dynamic roll tilt compensation, absence of vehicle-specific hazard weighting, and the methodological divide between physical sensors and simulation models---form the explicit research rationale for the CycleSafe architecture presented in subsequent chapters.
